\documentclass[11pt]{article}

\usepackage[a4paper,margin=28mm]{geometry}
\usepackage{amsmath,amssymb,amsthm}
\usepackage[colorlinks=true,urlcolor=blue,linkcolor=blue,citecolor=blue]{hyperref}

\newtheorem{theorem}{Theorem}
\newtheorem{lemma}[theorem]{Lemma}
\newcommand{\Ree}{\operatorname{Re}}
\newcommand{\NF}{\operatorname{N}}

\title{The Quasi-Riemann Hypothesis: A Zero-Free Half-Plane
$\Ree(s)>0.8749570698$}
\author{}
\date{8 October 2026}

\begin{document}
\maketitle

\begin{abstract}
We examine the numerical endpoint in the compensated-probe proof of the
seven-eighths zero-free half-plane in \cite{OAI}. We increase the total
prime length and shift the two physical scales in opposite directions.
An exact polynomial certificate bounds the resulting high-side exponent.
The low comparison includes the same shift in its principal signal.
These estimates give the zero-free boundary
$4374785349/5000000000=0.8749570698$.
The accompanying Lean files prove the moment input and continuation
argument against the upstream formalization. The full proof passes its
axiom audit and Comparator's kernel replay.
\end{abstract}

\section{Statement and geometry}

Let $F=\mathbb Q(\sqrt{-3})$. As in \cite[Equation (1.1b)]{OAI}, let
$\beta_*$ be the supremum of $1/2$ and the real parts of zeros of primitive
finite-order Hecke $L$-functions over $F$ in the critical strip. We use the
seven-eighths theorem of \cite{OAI}, so $\beta_*\le 7/8$. We also use the
arithmetic identities, detector estimates, reflected-energy lemma, additive
Gram bound, and continuation criterion proved there. The result therefore
depends on the correctness of \cite{OAI}. This note changes only the
numerical geometry and the estimates that the change affects.
Section~\ref{sec:low} checks the changed low comparison.
Section~\ref{sec:moments} supplies the perturbed source moments with
zero moment excess, and Section~\ref{sec:high} checks the high comparison
and analytic domains.

\begin{theorem}\label{thm:main}
Every finite-order Hecke $L$-function over $F$ and every Dirichlet
$L$-function, including the Riemann zeta function, has no zero in
\[
    \Ree s> B:=\frac{4374785349}{5000000000}
                =\frac78-\frac{214651}{5000000000}.
\]
The usual pole at $s=1$ for a principal character is allowed.
\end{theorem}

Put
\begin{equation}\label{eq:geometry}
\begin{gathered}
 \epsilon=\frac{214651}{5000000000},\qquad
 \kappa=\frac{12457523527}{15625000000000},\\
 t=\frac{42930200831}{250000000000000},\qquad
 \ell=\frac16+t,\qquad b=\frac18-2\kappa,\\
 l_x=\frac{17}{48}+\kappa-\frac t2,\qquad
 l_y=\frac{23}{48}-\kappa-\frac t2,\\
 h=\frac{13}{16}-\kappa+\frac{3t}{2},\qquad
 M=l_x+l_y=\frac56-t.
\end{gathered}
\end{equation}
These lengths satisfy $l_y-l_x=b$, $M+\ell=1$, and
$h=1-l_x+\ell$. The general Mellin-exponent identity in
\cite[Part II, compensated low estimate]{OAI} gives
\begin{equation}\label{eq:signal}
 C(s)=\frac{l_x}{2}+s-1+\frac h6
     =s-\frac{11}{16}+\frac\kappa3.
\end{equation}
Choose finitely many distinct positive slot lengths $\ell_j$ with
$\sum_j\ell_j=\ell$. Their maximum can be made smaller than any prescribed
positive mesh. The prime annuli at scales $Z^{\ell_j}$ are then pairwise
disjoint for all sufficiently large $Z$. The finite physical
probe and its exact Poisson identity in \cite{OAI} make sense with the
lengths in \eqref{eq:geometry}. All its arithmetic factors are unchanged.

\section{The low comparison}\label{sec:low}

\begin{lemma}\label{lem:low}
The normalized compensated probe of \cite{OAI}, formed with
\eqref{eq:geometry}, satisfies for every $\varepsilon>0$
\[
       |J_\eta(Z)|\ll_{\eta,\varepsilon}
             Z^{3/16+\kappa/3-t/4+\varepsilon}.
\]
\end{lemma}

\begin{proof}
Fix a rescaled subset of slots of total logarithmic length
$0\le d\le\ell$. In the notation of the compensated low estimate of
\cite{OAI}, its row and surviving-slot lengths are
\[
 M'=M-2d,\qquad \ell'=\ell-d.
\]
The proof of the reflected-row norm in \cite[Part II, compensated low estimate]{OAI}
first bounds one branch by $M'$ and the other by
$M'+(1+3\ell'-2M')/4$, up to arbitrary small losses. Its general
reflected-energy lemma permits these fixed bounded lengths. Substitution
gives the squared row norm
\begin{equation}\label{eq:row}
 \sum_{0<\NF m\ll Z^{M'}} |B^J_{m,\sigma}(Z)|^2
 \ll Z^{M'+((d-1/6+5t)/4)_++\varepsilon}.
\end{equation}
The additive Gram estimate is available throughout this range. Its scale
conditions follow from
\[
 l_x-\ell=\frac3{16}+\kappa-\frac{3t}{2}>0,\qquad
 l_y-\ell=\frac5{16}-\kappa-\frac{3t}{2}>0,
\]
\[
 M-2\ell=\frac12-3t>0,\qquad
 l_y-\ell-\frac{11b}{6}
   =\frac1{12}+\frac{8\kappa}{3}-\frac{3t}{2}>0.
\]
Thus its $P_a^2/Y'$ term is absorbed by $P_a^{1/6}$, exactly as in
\cite[Part II, additive Gram bound]{OAI}.
The compensated scale cap is now
\[
 2l_x-l_y=\frac{11}{48}+3\kappa-\frac t2,
 \qquad
 2l_x-l_y-\ell=\frac1{16}+3\kappa-\frac{3t}{2}>0.
\]
For a fixed calibration norm $q$ and an unselected tuple norm
$L\ll Z^\ell$, these inequalities give, eventually,
$qZ^{l_x}\le Z^{l_y}$ and
$LZ^{l_y}\le q^2Z^{2l_x}$.
They are the physical scale conditions used by the compensated Gram bound.
Cauchy--Schwarz, the count of
rescaled tuples, and their coefficients add to the basic low exponent
\[
 -d+\left(\frac{d-1/6+5t}{8}\right)_+.
\]
This is nonpositive for $0\le d\le\ell$. Indeed, its positive-part
argument is negative until $d>1/6-5t$, and thereafter
$d-1/6+5t<8d$. The empty subset has exponent zero. The direct bound
is consequently $Z^{l_x/2+b/12+\varepsilon}$.

The prime normalizer in \cite[Part II, principal normalizer]{OAI} is nonzero and has
inverse $Z^{o(1)}$ for every fixed positive slot length. Finally,
\[
 \frac{l_x}{2}+\frac b{12}
 =\frac3{16}+\frac\kappa3-\frac t4.
\]
This proves the normalized estimate.
\end{proof}

The signal at the proposed boundary is
\begin{equation}\label{eq:low-gap}
 \begin{gathered}
 C(B)=\frac3{16}+\frac\kappa3-\epsilon,\\
 C(B)-\left(\frac3{16}+\frac\kappa3-\frac t4\right)
 =\frac t4-\epsilon=\frac{831}{10^{15}}>0.
 \end{gathered}
\end{equation}
The formal continuation assembly uses low loss $10^{-13}$, which leaves
the strictly positive margin $731/10^{15}$.

\section{The perturbed source moments}\label{sec:moments}

The fourth-moment lemma with short prime factors in \cite{OAI} permits
moment parameter $\kappa_{\rm mom}=3/4$.
This parameter is distinct from the geometric skew $\kappa$ in
\eqref{eq:geometry}.
Its zero-free hypothesis is available because
$2\beta_*-1\le3/4$. The original detector application uses
$\kappa_{\rm mom}=3/4+2\Delta$ with $\Delta=\beta_*-7/8>0$.
Here the moment estimates use $\Delta=0$. The gap $\beta_*-B$, when
positive, is used separately to choose the detector losses.

We record the estimates needed for this application. Use the normalized
witness polynomials and physical slot factors of \cite{OAI}. In a source
fiber $\mathcal B$ at scale $U=Z^d$, write $r$ and $m$ for the inverse and
plain witness lengths, $I_{u,n}$ for the inverse witness with $n$
logarithmic derivatives, and $P_{u,j}$ for the plain witness with $j$
logarithmic derivatives. These include the real shifts in $[0,1]$ and
the height shifts in $[-H,H]$. The physical slot factors use
$\Ree z=17/50$ and $|\operatorname{Im}z|\le H$.
For a selected set $A$ of slots, let
$Q_A(u)$ be their physical product and put
\[
 \lambda_j=\ell_j/d,\qquad L_A=\sum_{j\in A}\lambda_j.
\]

\begin{lemma}\label{lem:moments}
Suppose $B<\beta_*\le7/8$. For every admissible choice of the perturbed
detector losses, let $\rho>0$ be its moment loss. There is a positive mesh
$\mu$ such that slot lengths with $\ell_j\le\mu/200$ satisfy the following
estimates. After fixing the slots and source, there is an exponent
$J\ge0$ independent of the target character. For each target there are
$C_\eta>0$ and a threshold such that the bounds hold for every source fiber with
$1/200\le d\le7/8$, $51/100<a\le7/8$, and retained rows
$Z^{1/100}\le\NF u\le Z^{d-\rho}$ above that threshold.
With $C_H=C_\eta(1+H)^J$, uniformly for $H\ge0$,
\begin{align}
 \sum_{u\in\mathcal B}|I_{u,n}Q_A(u)|^2
   &\le C_H U^{1+\rho},
   && n\le2,\quad r+2L_A<1,\quad 2r+8L_A<3,
   \label{eq:inverse-moment}\\
 \sum_{u\in\mathcal B}|P_{u,j}P_{u,k}Q_A(u)|^2
   &\le C_H U^{1+\rho},
   && j+k\le2,\quad 2m+\tfrac92 L_A\le1,
   \label{eq:plain-marked-moment}\\
 \sum_{u\in\mathcal B}|P_{u,j}P_{u,k}|^2
   &\le C_H U^{\max\{1,2m\}+\rho},
   && j+k\le2.
   \label{eq:plain-unmarked-moment}
\end{align}
The unmarked inverse raw-moment bounds also hold with the count parameters
chosen from the fixed source. Together these are the source moment bounds
with zero moment excess required by the detector.
\end{lemma}

\begin{proof}
Apply the certified energy theorem of \cite{OAI} with $\kappa_{\rm mom}=3/4$ and
stage error $\rho/4$. Its positive-width and zero-slot conclusions give a
positive fine mesh $\mu$. The condition $\ell_j\le\mu/200$ ensures
$\lambda_j\le\mu$ for every $d\ge1/200$. This does not restrict the total
slot length, since the number of distinct positive slots can be increased.

Fix the source's finite excluded prime set $S$, and let $M_S$ be its
product ideal. Use the fixed ideal $M_S\cap(72)$ in the energy source.
For each target there are only finitely many source labels. The norm of
the character modulus at every such label is at most $U^{\rho/4}$ once $Z$ is large
enough. The exceptional sixth-power-free rows are supported on a fixed
finite set of primes and have bounded norm. They are therefore absent
from the retained rows for large $Z$.

The source identities bound the sum of the squared plain witness products
by retained energy with the physical prime coefficients. The marked initial state
has width $1+\rho/4$. Its capacity condition follows from
$2m+6\kappa_{\rm mom}L_A\le1$, and the positive-width energy bound gives exponent
$1+\rho/2\le1+\rho$. For the unmarked pair, use initial width
$\max\{1,2m\}+\rho/4$ and the zero-slot energy bound. This gives
\eqref{eq:plain-marked-moment} and \eqref{eq:plain-unmarked-moment}.

For the inverse witness, the fixed-source inverse estimates of \cite{OAI}
apply under the two strict capacity inequalities in
\eqref{eq:inverse-moment}. Distinct slot lengths give disjoint prime
supports. Their fixed conductor deletions have modulus-one denominators
on these supports for large $Z$. The inverse raw-moment bound and the
marked inverse bound consequently apply to the perturbed slots.

The profile estimates bound all needed seminorms by a fixed power of
$1+H$. Taking a common exponent and summing the finitely many label
constants gives $J$ and $C_\eta$. The exponent is chosen before the
target. Thus the estimates hold at detector height $H=Z^{2\tau}$, and
the usual late choice of $\tau>0$ absorbs the polynomial height factor
into the prescribed detector loss.
\end{proof}

\section{The high comparison}\label{sec:high}

We first record the extension of the Euler regions used by the contour
arguments. The local factors below are those of
\cite[Equations (5.13a)--(5.13c)]{OAI}.

\begin{lemma}\label{lem:analytic}
The full compensated correction is holomorphic, with the bounds needed
for principal-row contour moves, on the second Euler region of \cite{OAI}
with its condition $\Ree s\ge7/8$ replaced by $\Ree s\ge B$. The
principal local quotient there satisfies
\[
       G_p/H_p=-1+O((\NF p)^{-B}).
\]
The small-row contour can use the first Euler region with
$\epsilon_0=37/100$.
\end{lemma}

\begin{proof}
The exact local defect formula in \cite[local Euler lemma]{OAI} bounds
$H_p-1$ in the original second
region by $O(Q^{-363/200})$ for $p\nmid u$ and $O(Q^{-33/40})$
for $p\mid u$, where $Q=\NF p$. Every local exponent involving
$\Ree s$ has coefficient of absolute value at most six. Moving its
lower boundary left by $\epsilon$ therefore gives the respective bounds
\[
 O(Q^{-363/200+6\epsilon})\quad\text{and}\quad
 O(Q^{-33/40+6\epsilon}).
\]
The first exponent is less than $-1$. The unramified product converges
normally, while the ramified product has the same $O_\varepsilon((\NF
u)^\varepsilon)$ bound. For $u=1$ its tail still tends to zero faster
than $P_0^{-4/5}$. Enlarging the fixed excluded set therefore gives the
same uniform nonvanishing of the principal product.

The four errors for the principal selected factor have exponents
$-\Ree s$, $-6\Ree z$, $4-5\Ree s-6\Ree z$, and
$1-\Ree w-6\Ree z$. On the extended region they are all at most
$-B$. The formula for $G_p$ has no division by $P_p$ or $H_p$, so the
full finite selected-tuple correction remains holomorphic at their zeros.
The bounds above and finite slot count give the all-height majorant used
in the principal contour proof of \cite{OAI}.

For small rows, $\beta_*>B$ implies
$\beta_*+e+1/2>1+37/100$. The first Euler region of \cite{OAI}
allows any positive $\epsilon_0$, so its unchanged local estimates apply
to the line $(\beta_*+e,1/2,17/50)$. The central bin contours already
lie in its region with $\epsilon_0=10e$. These are the only analytic
domain uses of the original $7/8$ lower bound in the shared contour
lemmas. Repeating their proofs with the extended region gives the same
principal, bin, and outer-row contour estimates.
\end{proof}

\begin{lemma}\label{lem:high}
If $B<\beta_*\le7/8$, then for
some $\sigma>0$ independent of the
target character,
\[
       |J_\eta(Z)-f_\eta(Z)|
          \ll_\eta Z^{C(\beta_*)-\sigma},
\]
where $f_\eta$ is the principal Mellin signal from \cite{OAI} with
\eqref{eq:signal}.
\end{lemma}

\begin{proof}
Lemma~\ref{lem:moments} supplies the fourth-moment estimate at
$\kappa_{\rm mom}=3/4$, because its
zero-free input is $\beta_*\le(1+\kappa_{\rm mom})/2=7/8$. Every detected bin has
$\delta=2a-1\le2\beta_*-1\le3/4$. Its plain prime capacity is thus
the baseline capacity in the proof of \cite[Part II, capacity
comparison]{OAI}, with no positive $\Delta/4$ surcharge.
The row exponent is bounded by $R_*(\delta,x)$ plus an arbitrarily small
loss, with $0\le x\le1/2$ and $q=x\delta$. The same selected-prime
argument applies since $\ell/h>8/39>7/37$ and
$5\ell-h>1/48$. A fixed sufficiently fine even slot mesh meets all
rounding conditions.

At row length $d=h$, the general high exponent relative to $C(B)$ is
\begin{equation}\label{eq:high-exponent}
 E_B(h)=a-B+h\left(\frac{17}{50}-\frac16\right)
       -a l_y-\frac\ell2+q\ell
       +h\left(R+\frac\delta2-\frac{17}{50}\right).
\end{equation}
Using $\delta=2a-1$ and holding $a,\delta,q,R$ fixed while changing
\eqref{eq:geometry} gives
\begin{equation}\label{eq:perturbation}
 E_B(h)-E_{7/8}(13/16)
 =\epsilon+t\left(-\frac12+\delta+q+\frac{3R}{2}\right)
       +\kappa\left(\frac23-R\right).
\end{equation}
We evaluate the combined exponent using the exact balanced row count
from \cite[Part II, compensated endpoint certificate]{OAI}. Put
$y=1/2-x$ and define
\[
 D(y)=\frac{37+34y}{18},\qquad
 P(y)=\frac{7+18y+8y^2}{9},\qquad
 J(\delta,y)=(5/6-\delta)D(y)+\delta P(y).
\]
Then $J\ge35/54>0$ on the parameter box, and
\[
 R_*=1-\delta+\frac{(5/6-\delta)\delta P(y)}{2J(\delta,y)}.
\]
With $R=R_*$ and the value of $t$ in \eqref{eq:geometry}, direct substitution gives
\[
 \begin{gathered}
 162000000000000000J(\delta,y)
   \left(-E_B(h)-\frac1{1250000000000}\right)
 =H(\delta,y),\\
 H(\delta,y):=A(y)\delta^2+L(y)\delta+K(y),
 \end{gathered}
\]
where
\begin{align*}
 A(y)&=38381036152738950+98112671277888624y\\
     &\quad+70349424015981576y^2+24024727795678656y^3,\\
 L(y)&=-29828491018092129-47206874749477924y\\
     &\quad-3166469409426772y^2,\\
 K(y)&=5795432885635430+5325532921935260y.
\end{align*}
The following identity certifies $H\ge0$ for every $y\ge0$:
\begin{align*}
 4A(y)H(\delta,y)
 &=(2A(y)\delta+L(y))^2+979807389619512241359\\
 &\quad+275799814009247843319007787580888y\\
 &\quad+1303439473219357232720782491774728y^2\\
 &\quad+1756569235069508453488622075444704y^3\\
 &\quad+501751386744843713686464752902256y^4.
\end{align*}
Every coefficient in the remainder is positive, as is $A(y)$. Hence
\begin{equation}\label{eq:high-gap}
 -E_B(h)\ge\frac1{1250000000000}.
\end{equation}
For a row-count loss $0\le\rho\le1/32$, put $R=R_*+\rho$.
The source bounds imply
$0\le R\le3/2$ and $0\le R+\delta/2-17/50\le2$.
For $\zeta,\mu\ge0$, frequency length $v\le h+\zeta$, and row scale $d$
with $d-v\le\mu$, the mixed exponent differs from its balanced endpoint by
\[
 h\rho+(v-h)(R+\delta/2-17/50)+(d-v)R
 \le h\rho+2\zeta+\frac32\mu.
\]
These losses, the additional physical-scale cost
$t(3e+\mathrm{mesh})+6\kappa e$,
and the remaining detector losses can all be chosen below the margin
in \eqref{eq:high-gap}.

The floor and intermediate-row margins in \cite{OAI} are respectively
$7/1200$ and $49/14400$ at the original geometry. The floor shift costs
$t(121/40+3e+\mathrm{mesh})+\kappa(67/100+6e)$,
including the changed row threshold. The
parameter budget absorbs this cost and the boundary change
$7/8-B=\epsilon$. In particular,
\[
 \epsilon+\frac{121t}{40}+\frac{67\kappa}{100}
 <\min\left(\frac7{1200},\frac{49}{14400}\right).
\]
The intermediate estimate can also be checked directly.
Before the small detector losses, its bounds are
$d\le1/2$, $q\le\delta/2$, and $R\le76/75-2\delta/3$.
Its exponent relative to $C(B)$ is therefore at most
\[
 -\frac{529}{2400}+\frac{25\delta}{96}+\epsilon
 +t\left(\frac1{100}+\frac{3\delta}{4}\right)
 +\kappa\left(\frac{49}{150}+\frac\delta2\right)
 <-\frac3{125}
 \qquad(0\le\delta\le3/4).
\]
The expression increases with $\delta$, and substitution of $\delta=3/4$
gives $-2440350045024453/10^{17}$.
The small-row margin is $63/800$.
Relative to the shifted signal, its geometric cost is
\[
 \frac{51t}{100}+\frac{49\kappa}{150}<\frac{63}{800}.
\]
For large rows, the absolute contour can still be taken sufficiently far
right after choosing $\zeta$. Lemma~\ref{lem:analytic} supplies the
principal contour, whose margins $l_y/20$ and $h/600$ remain positive.
Its local approximation supplies a positive power saving in each fixed
slot. The prime normalizer costs only $Z^{o(1)}$.

All detector, mesh, real-contour, and normalizer losses can be chosen
below these fixed positive margins before the target character is fixed.
The explicit central budget takes a common small parameter at most
$10^{-16}$ and bounds the later detector allowances by
$\mathrm{small}/(N+2000)$, where $N$ is the fixed slot count.
It reserves the high margin in \eqref{eq:high-gap}, including the row-count,
physical-scale, Mellin, and normalization costs.
The late height choice in \cite{OAI} then gives one $\sigma>0$ for every
target, as claimed.
\end{proof}

\begin{proof}[Proof of Theorem~\ref{thm:main}]
Suppose $\beta_*>B$. Lemma~\ref{lem:analytic} supplies a principal
factor holomorphic and bounded away from zero on $\Ree s>B$.
Lemma~\ref{lem:low} and \eqref{eq:low-gap} give the direct bound required
by the continuation criterion of \cite[Section 2]{OAI}. Its parameter
$\omega$ can be any sufficiently small positive number below
$\beta_*-B$. Lemma~\ref{lem:high} gives the other required bound with a
target-independent saving. The criterion contradicts $\beta_*>B$.
Hence $\beta_*\le B$. The transfer in \cite{OAI} from primitive Hecke
characters to imprimitive Hecke and Dirichlet characters proves the
stated assertions.
\end{proof}

\section*{Verification scope}

The rational inequalities in \eqref{eq:low-gap} and \eqref{eq:high-gap}
can be checked independently by the accompanying
\texttt{scripts/verify\_bounds.py}.
The accompanying files in \texttt{lean/} contain proofs of the geometry,
reflected-exponent calculation, high-side endpoint, normalized compensated
low estimate, local Euler bounds, and principal and global Euler product
extensions against the upstream OAI formalization.
The finite source multiplier bounds and
the principal source's $w$ and $z$ contour shifts, ordered double shift,
uniform contour tails, and $x$ contour transport also have Lean proofs at
the new boundary. The selected high-row local factors, weighted prime-slot
sum, physical-row dyad, Mellin integral, and changed small- and large-row
tail bounds, central cube truncation error, canonical ray-cube choice, and
normalized floor-row collection have proof files there as well.
The principal slot-ratio and normalized residue proofs use the
new boundary. The physical principal remainder at the changed lengths is
bounded on a fixed contour just right of $7/8$, leaving a power saving of
$1/4000$ relative to $\beta_*$ in the proposed interval. The contour at
$\beta_*+e$ used in the final comparison retains a saving of $1/3000$.
The perturbed source moment transport is implemented in
\texttt{OAIHighMomentTransport.lean}. It supplies the chosen fourth-moment
input used by the central-bin estimate and detector assembly.
\texttt{OAIHighUnconditional.lean} derives the prior $\beta_*\le7/8$ bound
from the upstream unconditional module and concludes the stronger Hecke,
Dirichlet, and zeta nonvanishing statements without a moment hypothesis.
The final theorem's axiom audit passes with only the standard Lean axioms
\texttt{propext}, \texttt{Classical.choice}, and \texttt{Quot.sound}.
It rejects any other transitive axiom dependency, including incomplete proofs.
The verification script requires the pinned upstream commit and a clean
upstream working tree.
Full Lean verification and Comparator replay of the updated bound pass.
The Comparator challenge uses upstream's Hecke definitions with the
bound changed to $B$. Comparator checks all three nonvanishing statements
and replays their proofs through the Lean kernel.
The proof uses Lean 4.34.1 and upstream revision
\texttt{adc7f1241b42e322a6451854ab7e4b4c146bf78a}.
Compiling this note checks the TeX. The separate Lean build checks the proof.
The stronger boundary in Theorem~\ref{thm:main} is an extension of
\cite{OAI}'s stated theorem.

This note, its Lean proofs, and the verification scripts were entirely
AI generated.

\begin{thebibliography}{1}
\bibitem{OAI}
OpenAI, \emph{The Quasi-Riemann Hypothesis: A Zero-Free Half-Plane
$\Ree(s)>7/8$}, 30 September 2026. Source consulted 7 October 2026.
\href{https://github.com/openai/math/blob/adc7f1241b42e322a6451854ab7e4b4c146bf78a/preprints/The-Quasi-Riemann-Hypothesis-September-30-2026/build/paper.tex}{Source TeX at the verified revision}.
\end{thebibliography}

\end{document}
